% =====================================================================
\section{G1 --- existence, and the node board}

\figpipeline

\figlattice
\label{ch:existence}
% =====================================================================

\subsection{What G1 asks}

G1 asks for a 4D continuum $\mathrm{SU}(N)$ gauge theory satisfying OS0--OS4. It is useful to
separate the statement into five nodes, because only one of them is open, and because that one is
the same species as the gap. The decomposition below is the corpus's
(\code{docs/PROOF-SKELETON.md} \S2, \code{research/unconditional/G1-ASSEMBLE-100.md}), with the
labels of this paper.

\begin{table}[t]
\centering
\footnotesize
\begin{tabularx}{\textwidth}{@{}l>{\raggedright\arraybackslash}Xc>{\raggedright\arraybackslash}p{3.2cm}@{}}
\toprule
\textbf{Node} & \textbf{Object} & \textbf{Local \%} & \textbf{Label}\\
\midrule
G1.0 & Fixed-volume assembly: L1 definition, L2 theorem, single-block $G''\ge\delta_W$,
  GL3 compactness, GL4 reflection-positivity propagation & $\sim95\%$ (all but L3) & \LT{} mod.\ L3\\
G1.1 & Perturbative/combinatorial kernel: BRST dimension-4 classes reduce to
  $\{\operatorname{Tr}F^2,\operatorname{Tr}F\tilde F\}$; asymptotic freedom absorbs the marginal
  log; $\sum_j g_j^4=\zeta(4)/(16\beta_0^4)<\infty$ & 100\% & \LT{} closed\\
G1.2 & \textbf{Uniform $Z_a$ $=$ L3}: $\sup_a Z_a^{-1}\|H_a\|_D<\infty$ $=$ L2 $+$ block
  Poincar\'e $+$ $\Lambda_{\mathrm{UV}}$ tower uniform & $\sim40\%$ & \LM{} open\\
G1.3 & $a\to0$ continuum via OS reconstruction, given L3 & $\sim70\%$ & \LC{} on L3\\
G1.4 & $L\to\infty$, OS2 (continuous O(4)), unique vacuum OS4 (clustering) & 0\% &
  \LM{} open, entangled with G2\\
\bottomrule
\end{tabularx}
\caption{G1 node board: existence. Percentages are node-local, not a percentage of Clay.}
\label{tab:g1}
\end{table}

\subsection{G1.1 --- the perturbative kernel is closed}

The perturbative kernel is worth stating in detail because it is a genuine closure and because it
shows what does \emph{not} go wrong. On the lattice, a dimension-four BRST-invariant local
operator can be built from the field strength in only two ways at the level of the classical
action: $\operatorname{Tr}F^2$ and $\operatorname{Tr}F\tilde F$. Any other candidate is either
BRST-exact (hence vanishes on physical states), a total derivative (hence a boundary term), or of
dimension $>4$ (hence irrelevant). The marginal logarithm associated with $\operatorname{Tr}F^2$
is absorbed by asymptotic freedom: the beta function is negative, so the marginal coupling flows
to zero in the ultraviolet and the log is not an obstruction to a uniform bound. The
combinatorial sum over perturbative orders is controlled by $\sum_j g_j^4<\infty$, a convergent
series at the Gaussian fixed point. This is the content of G1.1, and it is complete.

The closure matters for a negative reason: it rules out the hope that the ultraviolet obstruction
is a marginal-operator problem. The obstruction is not in the operator classification. It is in
the \emph{uniformity} of a force-variance bound (L3).

\subsection{G1.2 --- L3, the one open UV node}
\label{sec:l3}

L3 is the statement
\begin{equation}
\sup_a\ Z_a^{-1}\,\|H_a\|_{\mathcal D}\ <\ \infty ,
\label{eq:l3}
\end{equation}
where $H_a$ is the force Hessian of the lattice action and $\|\cdot\|_{\mathcal D}$ is the
Dirichlet-form norm on the space of gauge-field fluctuations, $Z_a$ the field-strength
renormalization constant. Equivalently, it is a cutoff-uniform bound on the force variance in
the ultraviolet tower. Three routes have been built towards it
(\code{G1-L3-TOWER-POINCARE.md}, \code{G1-L3-UV-FORCE-BOUND.md}, \code{G1-L3-VIA-FLOW.md}):
a tower Poincar\'e inequality, a direct UV force bound, and a gradient-flow route. None closes
it; each reduces it to the same single-sided variance bound, which is the UV face of the one
species (Chapter~\ref{ch:immune}).

\textbf{Why L3 is not the gap.} L3 is purely ultraviolet. Proving it would make fixed-volume
existence unconditional (G1.0, G1.3) but would not give a lower bound on the physical mass. The
two are orthogonal: L3 controls the short-distance fluctuations, the gap controls the long-
distance spectrum. The reason they are nevertheless ``one species'' is structural
(\S\ref{sec:uv-face}): both are one-sided upper bounds on a conditioned variance, written at
different scales. Building one tool that proves the inequality at both scales would close both
faces simultaneously --- a conjecture, not a theorem.

\textbf{Honest percentage.} The corpus assigns L3 $\sim40\%$. That number means: the statement is
reduced to a single variance inequality, the supporting analytic framework (block decomposition,
Poincar\'e on the UV tower, compactness) is in place, and what is missing is the inequality
itself. It is not 40\% of a proof; it is the fraction of the surrounding architecture that is
built around a missing strut.

\subsection{G1.3 --- from L3 to the continuum}

Given L3, the continuum limit is assembled as follows. The fixed-volume measure $\mu_a$ is
defined by the Wilson action and the Haar measure; L2 gives a uniform bound on the action
density; L3 gives the uniform renormalization constant; block Poincar\'e gives tightness of the
family $\{\mu_a\}$; and OS reconstruction then produces a measure $\mu$ on $\mathbb{R}^4$
satisfying the growth/analyticity axiom OS0, the invariance axiom OS1, and reflection
positivity.\footnote{The numbering follows the corpus, in which reflection positivity is OS3 and
the continuous O(4) action is OS2; in the conventional Osterwalder--Schrader list reflection
positivity is OS2. The labels are used only to track which axioms are proved, and the ambiguity
is recorded here so that the reader is not misled.} The one axiom that is not purely ultraviolet
is OS2 (continuity of the O(4) action and the existence of the infinite-volume limit), which is
entangled with G1.4.

The corpus's label is \LC{} on L3, at $\sim70\%$: the assembly is complete modulo the missing
inequality. This is the precise sense in which ``G1 is 100\% assembled modulo $\{L3,G2\}$'':
every input row is either a theorem, a reduction to L3, a reduction to the gap, or the same
species; there is no genuinely independent third open item. The master document records a
coarser figure of ``G1 $\approx70\%$'' (\code{docs/YANG-MILLS-MASTER.md} \S9); that number is a
summary of ``perturbative kernel closed, L3 open,'' not a separate claim. We report both and do
not average them.

\subsection{G1.4 --- infinite volume, OS2, and the downstream role of OS4}
\label{sec:g14}

The infinite-volume limit requires two things: OS2, the existence of a continuous O(4) action
and the infinite-volume limit of the Schwinger functions, and OS4, the uniqueness of the vacuum
(clustering). OS4 is \emph{downstream of the gap}: a uniform gap gives exponential clustering
for free, so a proof of G2 would supply OS4 (with an explicit rate). OS2 is partly independent:
it requires tightness and a uniform regularity estimate, which the campaign has in the
strong-coupling regime and not at the physical coupling. The corpus's assessment
(\code{G1-OS2-INFVOL.md}, \code{G1-UNIQUENESS-VACUUM.md}) is that OS2 is not a separate wall
but a corollary of the same uniformity that L3 and H1 need. This is why the node board lists
G1.4 at 0\% but does not count it as an independent open item: it is entangled with G2 by
construction, not by accident.

\subsection{The UV face of the one species}
\label{sec:uv-face}

The one-species identification is worth stating at this point in the paper, because it is the
bridge between Parts I and II. In the exact identity
\begin{equation}
G''(0)\;=\;E\bigl[H_{vv}\mid0\bigr]-\operatorname{Var}\bigl(\langle v,\nabla S\rangle\mid0\bigr)
+W_{\mathrm{full}},
\label{eq:species-identity}
\end{equation}
the variance enters with a minus sign. The gap is $G''(0)\ge\lambda$, so \emph{any} upper bound
on the conditioned variance improves the lower bound on $G''$. The same shape appears in the
ultraviolet: L3 is an upper bound on the force variance of the UV tower. The two faces differ in
scale, not in form. A tool that upper-bounds a conditioned variance at one scale upper-bounds it
at the other, and the campaign states that building such a tool would close L3 and H1
simultaneously, hence Clay. This is the ``one species, two scales'' reduction. The unification is
the conjecture; the dictionaries --- the exact Ward identity $E[a\,f]=1$, the residue--level-gap
lemma, and the identity above --- are theorems.

\figscales

\subsection{The dimensional correction}
\label{sec:dim}

One correction must be recorded because it is easy to get wrong and it changes what
``$G''\ge\lambda$'' means (\code{docs/PROOF-SKELETON.md}, G1-CONTINUUM-SCALING). For the
dimensionless Wilson action, the Hessian satisfies
\begin{equation}
G''_a=a^2\,m_{\mathrm{phys}}^2,\qquad
m_{\mathrm{lat}}=\sqrt{G''_a}=a\,m_{\mathrm{phys}},\qquad
\Delta=\lim m_{\mathrm{lat}}/a .
\label{eq:dim}
\end{equation}
Therefore a claim of the form ``at fixed $a$, $G''_a\ge\lambda$ with $\lambda$ a positive
\emph{constant}'' is dimensionally wrong for the gap: it would give
$m_{\mathrm{phys}}=\lim\sqrt{\lambda}/a=+\infty$, the free-box UV scale, a lattice artifact rather
than a particle. The correct target is a \textbf{ratio} $G''_a/a^2\ge\Delta_0^2$ at the
asymptotically-free scaling rate, i.e.\ $\lambda_a\sim a^2\Delta_0^2$, with $\lambda_a$
transported multiplicatively by $Z_O^2$. The transport itself is free (it is a square); what is
not free is the non-zero seed it multiplies. This distinction pervades the honest reading of the
numerics in Chapter~\ref{ch:numerics} and of the verdict in Chapter~\ref{sec:verdict}.

\textbf{A second dimension trap.} The unconditioned expected Hessian
$E[\operatorname{Hess}V]=E[\nabla V\otimes\nabla V]\succeq0$ is provably positive semidefinite for
every measure (Stein/IBP). So the gap cannot come from the \emph{expected} curvature; it comes
from the \emph{conditioning correction} --- the difference between the conditioned and
unconditioned fibre Hessians. The sharpest form of H1 (Appendix~\ref{app:sharp}) is stated
directly in this conditioned language, and the numerical audit of
Appendix~\ref{app:postaudit} shows that the unconditioned budget $E[f^2]\approx2$ is
ultraviolet/equipartition-pinned and cannot stand in for the scaling $\lambda_a\sim a^2\Delta_0^2$.

\subsection{The node board, in one table}

\begin{table}[t]
\centering
\footnotesize
\begin{tabularx}{\textwidth}{@{}ll>{\raggedright\arraybackslash}p{4.5cm}
  >{\raggedright\arraybackslash}X@{}}
\toprule
\textbf{Half} & \textbf{Object} & \textbf{Local completion} & \textbf{Status}\\
\midrule
G1 & 4D continuum existence & 100\% assembly modulo $\{L3,G2\}$ & \LT{}+\LC, one UV node open\\
G1.2 & UV uniform $Z_a$ $=$ L3 & $\sim40\%$ & \LM\\
G2 & Gap: region coverage & $\sim90\%$ & theorems G2.1--4\\
G2.5--6 & Gap: sign / proof & \textbf{0\%} & \LM{} $=$ H1, the species\\
Channel & $0^{++}$ identification $+$ no dissolution & 100\% & \LC\\
\textbf{Clay} & existence $\wedge$ gap & \textbf{0\%} & \textbf{open}\\
\bottomrule
\end{tabularx}
\caption{The node board in one table (see also Figure~\ref{fig:bars}).}
\label{tab:nodes}
\end{table}

The percentages are \textbf{node-local} --- progress toward each node's own goal --- and are
emphatically \textbf{not} a percentage proof of Clay. Clay is 0\% because the species is
unbuilt. The next chapter turns to G2, where the region coverage is high and the sign is zero.

\figbars
